% Standalone formula guide: compile this file as the main document.
% Sources: contents/introduction.tex and contents/chapter2_shortened.tex.
% Chapters 3--6 are reserved for subsequent batches; they are not included here.
% Bibliography is embedded below, so no separate .bib upload is required.
\documentclass[11pt,a4paper]{article}
\usepackage[T1]{fontenc}
\usepackage[utf8]{inputenc}
\usepackage[margin=25mm]{geometry}
\usepackage{amsmath,amssymb}
\usepackage[hidelinks]{hyperref}
\numberwithin{equation}{section}
\setlength{\parindent}{0pt}
\setlength{\parskip}{0.5em}
\allowdisplaybreaks
\title{Formula Guide: Chapters 1 and 2}
\author{}
\date{}
\begin{document}
\maketitle

\textbf{Scope.} This guide follows the introduction and the shortened background
chapter in their reading order. It covers every displayed equation and the
inline formulas, dimensions and ranges used there. Repeated symbols are
explained with their formula rather than listed as separate formulas.

\textbf{How to read the sources.} A citation identifies a supplied paper that
actually presents the equation or supports the stated mechanism. It does not
automatically identify the equation's original inventor. Standard definitions,
algebraic counts and project-specific adaptations are marked explicitly where
no originating paper has been verified in the supplied collection. Paper page
numbers refer to the supplied PDF, whose preprint pagination may differ from
the final publication. Equations added to express prose mathematically are
identified as explanatory expansions.

\textbf{Notation.} Symbols are local to each topic: $v$ is vertical displacement
in optical flow but variance in SimAM; $\alpha$ is magnification gain in EVM but
a class weight in focal loss; $\epsilon$ may mean a numerical stabiliser or a
label-smoothing strength. These quantities are not interchangeable.

\tableofcontents
\clearpage

\section{Chapter 1: Introduction}
\label{fg:chapter1}

\textbf{Source file:} \texttt{contents/introduction.tex}. No shortened
introduction exists in this checkout. The chapter contains \textbf{no explicit
inline or displayed mathematical formulas}. It describes the four binary
component switches, twelve valid configurations, 25 subject folds and the
156-clip working set in prose. It also introduces matched comparisons without
writing their calculation as an equation. There are therefore no Chapter 1
equations or equation-specific paper attributions to extract. Formal equations
from later chapters have not been inserted into this section.

\section{Chapter 2: Background}
\label{fg:chapter2}

\textbf{Source file:} \texttt{contents/chapter2\_shortened.tex}. Its phenomenon
section (2.1) contains definitions but no mathematical formulas. The formulas
begin with recording and input preparation. Chapter 2 contains sixteen displayed
equation blocks; the guide also explains its inline calculations and notation.

\subsection{Recording and input preparation: thesis section 2.2}

\subsubsection{Duration and the number of frame intervals}
\label{fg:capture}
\textbf{Thesis location:} 2.2.1, High-speed capture.

The chapter's inline expression is $0.3f_s$ for an event lasting $0.3$ seconds.
Its general form, written out here for explanation, is
\begin{equation}
    n_{\mathrm{intervals}}\approx \tau f_s.
    \label{fg:eq:capture}
\end{equation}
Here $\tau$ is duration in seconds and $f_s$ is frames per second. Multiplying
them gives a number of temporal intervals: $0.3\times200=60$, versus 7.5--9 at
25--30 fps. Exact recorded frame counts depend on endpoint alignment; intervals
and frames are not the same count. Subsampling also changes the spacing between
the images that reach the model.

\textbf{Source.} This is the definition of sampling rate, not a new equation
from a research paper. Yan et al.~\cite{yan2014}, apparatus description on PDF
pages 2--3, support the CASME II recording rate of 200 fps.

\subsubsection{Image dimensions}
\label{fg:geometry}
\textbf{Thesis locations:} 2.2.2--2.2.3.

The expressions $640\times480$ and $224\times224$ denote image width and height
in pixels. They are dimensions, not learned transformations. The recorded study
uses the original full frames and resizes them to a square; this changes aspect
ratio and does not perform face alignment or cropping.

\textbf{Source.} Yan et al.~\cite{yan2014}, PDF page 2, report the original
capture dimensions. The square resize and choice of raw frames are project
settings, verified against the working manifest and frame loader; they are not
implied by the database paper.

\subsection{Motion magnification: thesis section 2.3}

\subsubsection{Small-motion approximation}
\label{fg:evm}
\textbf{Thesis location:} 2.3.1, The Eulerian idea.

% SOURCE-DISPLAY-01
\begin{equation}
    I(x,t)=f(x+\delta(t))
       \approx f(x)+\delta(t)\frac{\partial f}{\partial x}.
    \label{fg:eq:evm-taylor}
\end{equation}
$I(x,t)$ is observed intensity at image position $x$ and time $t$; $f$ is the
reference image profile; $\delta(t)$ is a small displacement. The derivative
$\partial f/\partial x$ measures how rapidly intensity changes across position.
The equation says that moving a textured profile slightly changes its intensity
approximately in proportion to displacement. It drops higher-order terms, so
large displacements need not satisfy the approximation.

\textbf{Source.} Bai et al.~\cite{bai2021}, PDF page 2, section 3.1,
equation (2), present this first-order Taylor expansion. They explain an existing
Eulerian magnification formulation; this citation does not credit them with
inventing Taylor expansion or EVM.

\subsubsection{The amplified intensity signal}
% SOURCE-DISPLAY-02
\begin{equation}
    \widetilde I(x,t)\approx
       f(x)+(1+\alpha)\delta(t)\frac{\partial f}{\partial x}.
    \label{fg:eq:evm-gain}
\end{equation}
$\widetilde I$ is the reconstructed intensity and $\alpha$ is the gain applied
to the retained temporal variation. The factor $1+\alpha$ includes the original
motion term plus its amplified addition. Within the small-motion approximation,
this resembles displacement increased by that factor. It does not guarantee
that all amplified variation is useful facial motion.

\textbf{Source.} Bai et al.~\cite{bai2021}, PDF page 2, equations (3)--(4).
The thesis retains this rationale but implements a hard Fourier mask and shared
gain across Laplacian detail bands, rather than reproducing every filtering
choice described in that paper.

\subsubsection{Pyramid blur kernel and output range}
\label{fg:pyramid}
\textbf{Thesis locations:} 2.3.2 and 2.3.5.

The inline kernel is
\begin{equation}
    h=[1,4,6,4,1]/16.
    \label{fg:eq:blur}
\end{equation}
It is a five-tap weighted average: the centre receives the greatest weight and
the weights sum to one. Applying it separately along image rows and columns
smooths spatial detail before downsampling. The formula above is the blur
kernel; it is not a statement that every reconstruction operation uses the
same scaling.

\label{fg:clip}
After magnification the implementation clips intensities to $[0,255]$ and
converts them to 8-bit greyscale. This range is a storage constraint, not the
range of a physical displacement or a trained probability.

\textbf{Source.} Project implementation:
\texttt{Stage1\_DataPipeline/evm\_magnifier.py}. Bai et al.~\cite{bai2021}
support the Laplacian-pyramid mechanism, but the exact kernel and output
conversion are identified here as code choices, not attributed to that paper.

\subsection{Optical flow and strain: thesis section 2.4}

\subsubsection{Brightness constancy}
\label{fg:flow}
\textbf{Thesis location:} 2.4.2.
% SOURCE-DISPLAY-03
\begin{equation}
    I(x,y,t)=I(x+u\Delta t,y+v\Delta t,t+\Delta t).
    \label{fg:eq:brightness}
\end{equation}
The same moving image point is assumed to retain its intensity between two
times. Here $u$ and $v$ are horizontal and vertical velocities, so multiplication
by the time interval $\Delta t$ gives displacement. In the discrete project
tensors, the same symbols instead describe pixels of displacement between
selected frames; they must not be interpreted as calibrated velocity per second.

\textbf{Source.} OFF-ApexNet~\cite{liong2019a}, supplied PDF page 5,
equation (4) and the definitions immediately below it. That paper's notation is
equivalent after substituting its displacement variables. Its particular
onset--apex estimator is not the complete sequence-processing method of this thesis.

\subsubsection{The optical-flow constraint}
% SOURCE-DISPLAY-04
\begin{equation}
    I_xu+I_yv+I_t=0.
    \label{fg:eq:flow-constraint}
\end{equation}
$I_x$ and $I_y$ are spatial intensity derivatives and $I_t$ is the temporal
derivative. Taylor expansion of brightness constancy gives this constraint.
It is one equation with two unknown motion components; an edge constrains
motion along its intensity gradient but not along the edge itself. A flow
estimator therefore needs additional neighbourhood assumptions. Brightness
changes and large motion can violate the approximation.

\textbf{Source.} OFF-ApexNet~\cite{liong2019a}, PDF page 5,
equations (5)--(7). The thesis's use of Farneb\"ack flow is a separate estimator
choice, also reported for micro-expression recognition by Zhao
et al.~\cite{zhaoS2021}; the constraint alone does not specify that estimator.

\subsubsection{The small-strain tensor}
\label{fg:strain}
\textbf{Thesis location:} 2.4.3.

Write the displacement vector as $\mathbf{u}=[u,v]^{\mathsf T}$. Its gradient
contains the spatial derivatives of both components; superscript $\mathsf T$
means transpose. The symmetric part is
% SOURCE-DISPLAY-05
\begin{equation}
\begin{aligned}
    \varepsilon&=\tfrac12\left(\nabla\mathbf{u}
                         +(\nabla\mathbf{u})^{\mathsf T}\right),\\
    \varepsilon_{xx}&=\frac{\partial u}{\partial x},\qquad
    \varepsilon_{yy}=\frac{\partial v}{\partial y}.
\end{aligned}
    \label{fg:eq:strain-normal}
\end{equation}
% SOURCE-DISPLAY-06
\begin{equation}
    \varepsilon_{xy}=\varepsilon_{yx}
       =\tfrac12\left(\frac{\partial u}{\partial y}
                    +\frac{\partial v}{\partial x}\right).
    \label{fg:eq:strain-shear}
\end{equation}
$\varepsilon_{xx}$ and $\varepsilon_{yy}$ describe local extension or
compression; the equal off-diagonal terms describe shear. A uniform translation
has constant displacement and therefore zero strain. This does not confer
immunity to arbitrary head rotation, perspective effects or flow errors.

\textbf{Source.} Shreve et al.~\cite{shreve2011}, PDF page 2,
section III-B, equations (2)--(3). Although the paper uses the term finite
strain, it explicitly assumes small motion. The symmetric-gradient expression
shown here is the linearised, small-strain approximation, not a general
finite-deformation tensor.

\subsubsection{Implemented strain magnitude and the Frobenius comparison}
\label{fg:magnitude}
% SOURCE-DISPLAY-07
\begin{equation}
    \varepsilon_{\mathrm{mag}}
       =\sqrt{\varepsilon_{xx}^2+\varepsilon_{yy}^2
                    +\varepsilon_{xy}^2}.
    \label{fg:eq:strain-magnitude}
\end{equation}
This nonnegative scalar combines the two normal components and counts shear
once. Squaring removes their signs, so deformation direction is lost. The
original flow components accompany it as separate channels.

For comparison, the chapter's prose notes the extra shear contribution in the
Frobenius norm. Written out as an explanatory expansion, it is
\begin{equation}
    \|\varepsilon\|_F
      =\sqrt{\varepsilon_{xx}^2+\varepsilon_{yy}^2+2\varepsilon_{xy}^2}.
    \label{fg:eq:frobenius}
\end{equation}
Both symmetric off-diagonal entries enter this norm. The difference changes
shear's relative weight and is not merely a single global rescaling.

\textbf{Source.} The tensor components come from Shreve
et al.~\cite{shreve2011}; the exact three-term scalar is verified in
\path{Stage1_DataPipeline/step2_extraction/flow_strain_extractor.py}.
It is not claimed to be Shreve's identical aggregation rule. The Frobenius
expression follows from the standard matrix-norm definition.

\subsubsection{Tensor shape and channel normalisation}
\label{fg:tensor}
\textbf{Thesis location:} 2.4.4.

The channel order $[u,v,\mathrm{strain}]$ produces shape $(3,32,224,224)$:
three channels, 32 adjacent-frame fields and a $224\times224$ spatial grid.
The 32 fields come from 33 selected frames. These are dimensions, not a
formula from a paper, and the spatial grid covers the resized full frame.

\label{fg:input-normalisation}
The chapter describes two transformations in prose. Their code-equivalent
formulas are added here for clarity:
\begin{equation}
    z_i=\frac{x_i-x_{\min}}{x_{\max}-x_{\min}+10^{-8}},
    \qquad
    \widehat a_i=\frac{a_i-\mu_a}{s_a+10^{-6}}.
    \label{fg:eq:input-normalisation}
\end{equation}
Extrema, mean and standard deviation are computed separately for each channel
over its $T\times H\times W$ values. Extraction maps into $[0,1]$.
In the second expression, $a$ is the stored channel after any training
augmentation; evaluation applies no augmentation. Loading standardises $a$,
using its own mean $\mu_a$ and standard deviation $s_a$. The loader's standard
deviation uses the default sample correction. The small constants prevent
division by zero, and a constant channel need not have unit variance. The
standardisation stabiliser is added to the standard deviation, unlike the
inside-square-root stabiliser in batch normalisation below.

\textbf{Source.} Project implementation in
\texttt{flow\_strain\_extractor.py} and \texttt{Ablation\_Study/dataset.py}.
These are standard normalisation operations with code-specific scope and
constants; no originating research paper is asserted for their exact combination.

\subsection{Neural-network fundamentals: thesis section 2.5}

\subsubsection{A computational unit}
\label{fg:unit}
\textbf{Thesis location:} 2.5.1.
\begin{equation}
    a=\phi(w^{\mathsf T}x+b).
    \label{fg:eq:unit}
\end{equation}
The input vector $x$ is combined with learned weights $w$ by a dot product;
the bias $b$ shifts that result; the activation $\phi$ produces output $a$.
The activation's shape is fixed while the surrounding weights are learned.

\textbf{Source.} This is the standard affine-unit definition used in the
background chapter, not an invention attributed to a particular supplied paper.
OFF-ApexNet~\cite{liong2019a}, PDF page 6, equation (12), provides a related
convolutional weighted-sum, bias and activation formulation; it is contextual
support rather than the literal origin of this compact vector notation.

\subsubsection{Rectified linear activation}
\label{fg:relu}
\textbf{Thesis location:} 2.5.2.
\begin{equation}
    \operatorname{ReLU}(x)=\max(0,x).
    \label{fg:eq:relu}
\end{equation}
Negative inputs become zero and positive inputs pass through unchanged. This
nonlinearity prevents a stack of affine layers from reducing to one affine
mapping. The source chapter names GELU as well, but gives no GELU equation.

\textbf{Source.} OFF-ApexNet~\cite{liong2019a}, PDF page 6, equation (13),
explicitly presents this formula. It is a supplied source for the definition,
not the original ReLU paper.

\subsubsection{Convolutional parameter count}
\label{fg:convparams}
\textbf{Thesis location:} 2.5.4.
% SOURCE-DISPLAY-08
\begin{equation}
    N_{\mathrm{params}}=C_{\mathrm{out}}
       (C_{\mathrm{in}}k_Hk_W+1).
    \label{fg:eq:convparams}
\end{equation}
For an ungrouped convolution, each of the $C_{\mathrm{out}}$ output channels has $C_{\mathrm{in}}k_Hk_W$
weights and one bias. Here $C_{\mathrm{in}}$ counts input channels and
$k_H,k_W$ are kernel height and width. Weight sharing makes this count
independent of image size, although computation and activation memory still
grow with spatial resolution. If bias is disabled, the $+1$ is omitted; the
actual stem's convolution layers disable it.

\textbf{Source.} Direct parameter counting, not a borrowed empirical result
or a formula needing an invented source citation. The bias setting is verified
in \texttt{Ablation\_Study/models.py}.

\subsubsection{Pooling notation}
\label{fg:pool}
\textbf{Thesis location:} 2.5.5.

The inline $2\times2$ denotes the spatial pooling window. With stride two and
compatible dimensions, it halves height and width. It is a shape setting;
the shortened chapter does not display a pooling equation. Neither maximum
nor average pooling learns weights. A later projection is a separate learned
operation. This interpretation follows the implemented layer definitions,
rather than attributing the exact settings to a paper.

\subsection{Network blocks: thesis section 2.6}

\subsubsection{Three-dimensional tensor and kernel notation}
\label{fg:conv3d}
\textbf{Thesis location:} 2.6.1.

The tensor axes $(B,C,T,H,W)$ mean batch, channels, time, height and width.
The kernel $(k_T,k_H,k_W)$ directly mixes neighbouring times only when
$k_T>1$. The stem uses $(1,3,3)$ kernels, $(0,1,1)$ padding and $(1,2,2)$
pooling. Thus these kernels preserve time length and perform spatial
convolution. Whole-clip statistics in normalisation or SimAM can still create
dependence across frames. All these tuples are implementation settings,
verified in \texttt{Ablation\_Study/models.py}.

\subsubsection{Batch normalisation}
\label{fg:bn}
\textbf{Thesis location:} 2.6.2.
% SOURCE-DISPLAY-09
\begin{equation}
    \widehat x=\frac{x-\mu_c}{\sqrt{\sigma_c^2+\epsilon}},
    \qquad y=\gamma_c\widehat x+\beta_c.
    \label{fg:eq:bn}
\end{equation}
For channel $c$, subtract mean $\mu_c$, divide by the stabilised standard
deviation and apply learned scale $\gamma_c$ and shift $\beta_c$. Training
statistics for BatchNorm3d span $(B,T,H,W)$; evaluation uses running estimates.
These activation statistics differ from per-sample input normalisation.

\textbf{Source.} A standard layer definition, consistent with the
\texttt{BatchNorm3d} layers in \texttt{Ablation\_Study/models.py}.
The source chapter gives no originating batch-normalisation paper. No such
paper is claimed here merely because another reviewed model uses the layer.

\subsubsection{SimAM statistics}
\label{fg:simamstats}
\textbf{Thesis location:} 2.6.3.
% SOURCE-DISPLAY-10
\begin{equation}
    \mu=\frac1M\sum_i x_i,\qquad
    v=\frac1{M-1}\sum_i(x_i-\mu)^2.
    \label{fg:eq:simamstats}
\end{equation}
$x_i$ is an activation at one position in a channel, $M$ counts positions,
$\mu$ is their mean and $v$ their variance estimate. Here $M=THW$, so each
sample's statistics span the full clip; the original image setting uses $HW$.
Both sums run over all positions $i=1,\ldots,M$, and the sample variance
requires $M>1$, as satisfied by this feature volume.

\textbf{Source.} Yang et al.~\cite{yang2021}, PDF pages 4--5,
equation (5) and Figure 3. The printed variance definition divides by $M$,
but Figure 3's reference code divides by $M-1$. The thesis follows the code
and extends its context from images to clips. It must not be described as
an unchanged copy of the printed variance equation.

\subsubsection{SimAM energy and multiplicative gate}
\label{fg:simamgate}
% SOURCE-DISPLAY-11
\begin{equation}
\begin{aligned}
    E_i^{-1}&=\frac{(x_i-\mu)^2}{4(v+\lambda)}+\frac12,\\
    \widetilde x_i&=x_i\operatorname{sigmoid}(E_i^{-1}).
\end{aligned}
    \label{fg:eq:simamgate}
\end{equation}
The squared deviation measures distinctiveness from the channel mean.
Dividing by $v+\lambda$ puts it in the context of channel variation;
$\lambda=10^{-4}$ prevents an unstable division for low variance.
$E_i^{-1}$ is reciprocal energy. The sigmoid turns it into a bounded gate
that scales $x_i$, preserving score order but compressing its values.
No learned gate parameters are added, although calculating the gate costs
time and memory. Distinctive activation is not automatically useful motion.

\textbf{Source.} Yang et al.~\cite{yang2021}, equations (5)--(6) and
Figure 3. Taking the reciprocal of the paper's minimum-energy expression
produces the displayed score; the variance convention and clip context are
the implementation adaptations described above.

\subsubsection{Scaled dot-product attention}
\label{fg:attention}
\textbf{Thesis location:} 2.6.4.
% SOURCE-DISPLAY-12
\begin{equation}
    \operatorname{Attention}(Q,K,V)
      =\operatorname{softmax}\left(\frac{QK^{\mathsf T}}{\sqrt{d_k}}\right)V.
    \label{fg:eq:attention}
\end{equation}
$Q,K,V$ are query, key and value matrices derived from input tokens.
$QK^{\mathsf T}$ scores each query against the keys, and $d_k$ is the key
dimension. Dividing by $\sqrt{d_k}$ controls the growth of dot products;
row-wise softmax produces nonnegative weights summing to one. Multiplication
by $V$ makes each output a weighted combination of values. This is the
single-head operation, not the complete transformer block.

\textbf{Source.} Dosovitskiy et al.~\cite{dosovitskiy2021}, supplied PDF
page 13, Appendix A, equations (5)--(8), explicitly define this attention
operation with per-head dimension $D_h$, equivalent to $d_k$ here. This is a
verified supplied presentation of the standard transformer operation, not
its original invention. The thesis uses fixed positions; the equation alone
contains no position information and does not establish temporal-order sensitivity.

\subsubsection{The learned no-CNN projection}
\label{fg:projection}
\textbf{Thesis location:} 2.6.5.
\begin{equation}
    N_{\mathrm{projection}}=48\times96+96=4{,}704.
    \label{fg:eq:projection}
\end{equation}
Pooling three channels to $4\times4$ gives 48 features per frame. The
$48\rightarrow96$ linear layer then learns 96 sets of 48 weights and 96
biases. Removing the CNN therefore does not remove all learned input
processing. This is direct arithmetic from \texttt{RawPatchEmbedding} in
\texttt{Ablation\_Study/models.py}, not a count imported from a paper.

\subsection{Training under scarcity and skew: thesis section 2.7}

\subsubsection{Cross-entropy for a hard target}
\label{fg:ce}
\textbf{Thesis location:} 2.7.1.
\begin{equation}
    L_{\mathrm{CE}}=-\log p_t.
    \label{fg:eq:ce}
\end{equation}
Here $t$ indexes the true class, not time, and $p_t$ is its predicted
probability. A confident correct prediction has small loss; a small true-class
probability has large loss. The logarithm is natural. Under ordinary sampling,
frequent classes contribute more examples, but class frequency alone does not
determine each example's gradient.

\textbf{Source.} Standard negative log-likelihood. Zhang
et al.~\cite{zhang2022}, supplied PDF page 6, section 3.4, equation (15),
present multiclass cross-entropy; reducing its class sum for a one-hot target
gives this per-example expression. Zhao et al.~\cite{zhaoS2021} present the
focal extension below, whose unweighted zero-exponent case also reduces to it.

\subsubsection{Focal loss}
\label{fg:focal}
\textbf{Thesis location:} 2.7.2.
% SOURCE-DISPLAY-13
\begin{equation}
    \operatorname{FL}(p_t)=-\alpha_t(1-p_t)^\gamma\log p_t.
    \label{fg:eq:focal}
\end{equation}
For $\gamma>0$, the factor $(1-p_t)^\gamma$ suppresses easy, confident examples
relative to hard ones. $\alpha_t$ is an optional class weight, serving a
different purpose from difficulty weighting. The study uses $\gamma=2$.
For example, the modulating factor is 0.01 at $p_t=0.9$ and 0.81 at $p_t=0.1$.
These are modulation factors, not complete loss values.

\textbf{Source.} Zhao et al.~\cite{zhaoS2021}, PDF page 7, section 3.2.2.2,
equation (10), apply focal loss to micro-expression recognition and report
$\gamma=2$. This is an application source, not a claim that Zhao et al.
originated focal loss. The thesis additionally applies label smoothing,
which must be distinguished from this unsmoothed expression.

\subsubsection{Inverse-frequency sampling}
\label{fg:sampling}
\textbf{Thesis locations:} 2.7.3--2.7.4.
\begin{equation}
    w_i=\frac1{n_{c_i}}.
    \label{fg:eq:sampling}
\end{equation}
$n_{c_i}$ is the number of training examples in example $i$'s class, counted
within the current fold. Each represented class has total weight
$n_c(1/n_c)=1$, so normalised sampling probabilities give those classes
equal mass. This is an expectation over replacement draws; individual
batches need not be balanced and an epoch can repeat or omit clips.
The held-out loader is not rebalanced. When this sampler is active, the
implementation disables loss-side class weights to avoid applying two
frequency corrections at once.

\textbf{Source.} The weights and replacement rule are project choices in
\texttt{tools/run\_ablation\_experiments.py}. The equal-mass result is their
algebraic consequence, not an independently sourced research finding.

\subsubsection{Smoothed targets and smoothed cross-entropy}
\label{fg:smoothing}
\textbf{Thesis location:} 2.7.5.
\begin{equation}
    q_c=(1-\epsilon)\mathbf{1}[c=t]+\epsilon/C.
    \label{fg:eq:smoothing-target}
\end{equation}
There are $C$ classes. The indicator is one for the true class and zero
otherwise; $\epsilon$ redistributes some target mass uniformly. With
$C=3$ and $\epsilon=0.05$, the true-class target is approximately 0.9667 and
each other target approximately 0.0167. The targets sum to one.
% SOURCE-DISPLAY-14
\begin{equation}
    L_{\mathrm{smooth}}
      =-(1-\epsilon)\log p_t
        -\frac{\epsilon}{C}\sum_{c=1}^{C}\log p_c.
    \label{fg:eq:smoothing-loss}
\end{equation}
This is cross-entropy against the smoothed targets, obtained by substituting
them into the sum over class log-probabilities. It discourages extreme
confidence rather than preferentially weighting a minority class.

The chapter also describes the implemented combination in prose. Written
out as an explanatory expansion, it is
\begin{equation}
    L_{\mathrm{implemented}}=\alpha_t(1-p_t)^\gamma L_{\mathrm{smooth}},
    \label{fg:eq:combined-loss}
\end{equation}
with $\alpha_t$ effectively one when class weighting is disabled. The focal
factor still uses the unsmoothed hard-target probability $p_t$. This is the
per-example loss; the implemented default reduction averages it over the batch.

\textbf{Source.} Exact combination and smoothing strength verified in
\texttt{Ablation\_Study/losses.py}. The target distribution is a standard
definition and the loss is its algebraic expansion. No supplied paper is
claimed to originate this exact focal-plus-smoothing implementation; Zhao
et al.~\cite{zhaoS2021} support the focal component only.

\textbf{Remaining training prose.} Section 2.7.6 describes horizontal flipping
and negating the horizontal-flow channel $u$; it gives no further displayed
formula. Section 2.7.7 names AdamW, warmup and cosine annealing but writes no
update or schedule equation. Their full equations have not been invented as
extractions from this chapter.

\subsection{Evaluation: thesis section 2.8}

\subsubsection{Precision, recall and class F1}
\label{fg:metrics}
\textbf{Thesis location:} 2.8.1.
% SOURCE-DISPLAY-15
\begin{equation}
\begin{aligned}
    P_c&=\frac{TP_c}{TP_c+FP_c},\qquad
    R_c=\frac{TP_c}{TP_c+FN_c},\\
    F1_c&=\frac{2TP_c}{2TP_c+FP_c+FN_c}.
\end{aligned}
    \label{fg:eq:metrics}
\end{equation}
For class $c$, $TP_c$ counts correct predictions of that class, $FP_c$ counts
other examples incorrectly predicted as that class, and $FN_c$ counts its
examples assigned elsewhere. Precision asks how many positive predictions
are correct; recall asks how many actual positives were found. F1 balances
both through their harmonic mean. In the confusion matrix, entry $(i,j)$
counts true class $i$ predicted as class $j$.

The project sets undefined class scores to zero. Changing a single prediction
does not have a fixed effect on F1 because both its numerator and denominator
depend on the existing confusion counts.

\textbf{Source.} OFF-ApexNet~\cite{liong2019a}, supplied PDF page 9,
equations (20)--(22), explicitly gives recall, precision and their harmonic-mean
F1; substituting the first two into the third gives the count-based F1 above.
See et al.~\cite{see2019}, PDF page 3, equation (3), give the direct count-based
F1 expression. These support the standard definitions, not a claim that those
papers originated precision or recall.

\subsubsection{Macro F1 and micro F1}
\label{fg:macro}
\textbf{Thesis location:} 2.8.2.
% SOURCE-DISPLAY-16
\begin{equation}
    F1_{\mathrm{macro}}=\frac1C\sum_{c=1}^{C}F1_c.
    \label{fg:eq:macro}
\end{equation}
Every class contributes equally, regardless of support. Thus Negative does
not receive greater explicit weight merely because it supplies more clips.
The chapter describes micro aggregation in prose; its explanatory expansion is
\begin{equation}
\begin{aligned}
    F1_{\mathrm{micro}}&=
    \frac{2\sum_c TP_c}{2\sum_c TP_c+\sum_c FP_c+\sum_c FN_c}\\
    &=\frac{N_{\mathrm{correct}}}{N}
    \quad\text{for single-label multiclass prediction}.
\end{aligned}
    \label{fg:eq:micro}
\end{equation}
Here $N$ is the number of evaluated clips and $N_{\mathrm{correct}}$ is the
number classified correctly. Each incorrect classification contributes one false positive and one false
negative across the classes. Micro F1 therefore equals accuracy in this task;
the identity does not hold for every possible classification setting.

\textbf{Source.} Macro F1 and the pooled class-average construction are
supported by See et al.~\cite{see2019}, PDF page 3. The micro-F1 equality is
an algebraic consequence of single-label prediction, not a new experimental
result from that paper.

\subsubsection{Pooling folds before calculating F1}
\label{fg:fold-pooling}
\textbf{Thesis locations:} 2.8.3--2.8.5.

The chapter states the following distinction in prose. Let $A^{(k)}$ be the
confusion matrix of held-out subject $k$ and let $K$ be the number of folds.
An explicit explanatory form is
\begin{equation}
\begin{aligned}
    A_{\mathrm{pooled}}&=\sum_{k=1}^{K}A^{(k)},\\
    F1_{\mathrm{macro}}(A_{\mathrm{pooled}})
       &\ne \frac1K\sum_{k=1}^{K}F1_{\mathrm{macro}}(A^{(k)})
       \quad\text{in general}.
\end{aligned}
    \label{fg:eq:pooling}
\end{equation}
Compute class F1 after summing counts for the primary pooled score; do not
substitute the average of fold scores. F1 is nonlinear and folds can lack
classes. Equality is possible in special cases, so the inequality is not
asserted for every dataset. Subject-disjoint folds also do not undo selection
bias when the same held-out labels select a checkpoint and score it.

\textbf{Source.} See et al.~\cite{see2019}, PDF page 3, explicitly accumulate
confusion counts over LOSO folds before computing Unweighted F1. Its printed
equation (4) appears to omit the summation; the guide follows its explanatory
prose and the mathematically complete class average. The matrix notation above
is a formalisation added by this guide. The thesis uses $K=25$; its stored
\texttt{macro\_f1} field is the mean of fold scores, whereas its reported
primary score is reconstructed from pooled counts.

\subsection{Coverage boundary}
All sixteen displayed equation blocks in the shortened Chapter 2 are included
above, together with its inline formulas, tuple notation and numerical ranges.
The concluding section 2.9 has no further formulas. Topics mentioned without
an equation have not been expanded into a separate textbook treatment.
Chapters 3--6 will be separate subsequent additions.

\begin{thebibliography}{9}
\bibitem{yan2014}
W.-J. Yan et al., ``CASME II: An Improved Spontaneous Micro-Expression Database
and the Baseline Evaluation,'' \emph{PLoS ONE}, vol. 9, no. 1, e86041, 2014.

\bibitem{bai2021}
Mengjiong Bai, Roland Goecke and Damith Herath, ``Micro-Expression Recognition
Based on Video Motion Magnification and Pre-Trained Neural Network,''
\emph{IEEE International Conference on Image Processing}, pp. 549--553, 2021.

\bibitem{liong2019a}
Y. S. Gan, Sze-Teng Liong, Wei-Chuen Yau, Yen-Chang Huang and Lit Ken Tan,
``OFF-ApexNet on Micro-Expression Recognition System,''
\emph{Signal Processing: Image Communication}, vol. 74, pp. 129--139, 2019.
The supplied preprint lists Liong first; equation locators above refer to that PDF.

\bibitem{zhaoS2021}
Sirui Zhao et al., ``A Two-Stage 3D CNN Based Learning Method for Spontaneous
Micro-Expression Recognition,'' \emph{Neurocomputing}, vol. 448,
pp. 276--289, 2021.

\bibitem{shreve2011}
Matthew Shreve, Sridhar Godavarthy, Dmitry Goldgof and Sudeep Sarkar,
``Macro- and Micro-Expression Spotting in Long Videos Using Spatio-Temporal
Strain,'' \emph{IEEE International Conference on Automatic Face and Gesture
Recognition}, 2011.

\bibitem{yang2021}
Lingxiao Yang, Ru-Yuan Zhang, Lida Li and Xiaohua Xie,
``SimAM: A Simple, Parameter-Free Attention Module for Convolutional Neural
Networks,'' \emph{International Conference on Machine Learning}, 2021.

\bibitem{dosovitskiy2021}
Alexey Dosovitskiy et al., ``An Image Is Worth $16\times16$ Words:
Transformers for Image Recognition at Scale,''
\emph{International Conference on Learning Representations}, 2021.

\bibitem{zhang2022}
Liangfei Zhang, Xiaopeng Hong, Ognjen Arandjelovi{\'c} and Guoying Zhao,
``Short and Long Range Relation Based Spatio-Temporal Transformer for
Micro-Expression Recognition,'' \emph{IEEE Transactions on Affective
Computing}, vol. 13, no. 4, pp. 1973--1985, 2022.

\bibitem{see2019}
John See, Moi Hoon Yap, Jingting Li, Xiaopeng Hong and Su-Jing Wang,
``MEGC 2019---The Second Facial Micro-Expressions Grand Challenge,''
\emph{IEEE International Conference on Automatic Face and Gesture
Recognition}, 2019.
\end{thebibliography}
\end{document}
